\uuid{cF8j}
\titre{Composition de convolutions et champ réceptif}
\niveau{L3}
\module{Informatique / Signal}
\chapitre{Réseaux de neurones}
\sousChapitre{Empilement de couches convolutives}
\theme{Champ réceptif, comptage de paramètres, profondeur vs largeur}
\auteur{}
\datecreate{2026-03-23}
\organisation{}
\difficulte{3}

\contenu{
\texte{
  On empile deux couches de convolution successives, chacune avec un noyau $3 \times 3$,
  padding $= 0$ et stride $= 1$.
}
\begin{enumerate}
\item
  \question{Après la première couche, quelle est la taille de la sortie si l'entrée est $7 \times 7$ ?}
  \reponse{
    $\lfloor(7+0-3)/1\rfloor+1 = 5$. Sortie : $5 \times 5$.
  }

\item
  \question{Après la deuxième couche, quelle est la taille finale ?}
  \reponse{
    $\lfloor(5+0-3)/1\rfloor+1 = 3$. Sortie finale : $3 \times 3$.
  }

\item
  \question{Quel est le champ réceptif d'un pixel de la sortie finale (c'est-à-dire la zone
  de l'image d'entrée qui a influencé sa valeur) ? Montrer le raisonnement.}
  \reponse{
    Un pixel de la sortie de la couche 2 dépend d'une fenêtre $3 \times 3$ de la
    sortie de la couche 1.
    Chaque pixel de cette fenêtre $3 \times 3$ dépend lui-même d'une fenêtre $3 \times 3$
    de l'image d'entrée.
    La fenêtre totale couverte dans l'entrée est donc de taille
    $(3-1) + (3-1) + 1 = 5$, soit un champ réceptif de $\mathbf{5 \times 5}$.

    Formule générale pour $L$ couches de noyaux $k \times k$ (stride $= 1$, pas de padding) :
    \[
      \text{champ réceptif} = L(k-1) + 1.
    \]
    Ici : $2(3-1)+1 = 5$.
  }

\item
  \question{Comparer avec une unique convolution $5 \times 5$ appliquée directement sur l'image $7 \times 7$ (padding $= 0$). Même champ réceptif ? Même taille de sortie ?}
  \reponse{
    Taille de sortie : $\lfloor(7+0-5)/1\rfloor+1 = 3 \times 3$. Identique.

    Champ réceptif : $5 \times 5$. Identique.

    Les deux architectures voient la même zone de l'image pour chaque pixel de sortie.
  }

\item
  \question{Combien de paramètres comporte chaque approche (sans biais) ?
  Quelle configuration est plus économique ?}
  \reponse{
    \begin{itemize}
      \item \textbf{Deux couches $3 \times 3$} (1 seul filtre par couche, pour simplifier) :
        $3^2 + 3^2 = 9 + 9 = 18$ paramètres.
      \item \textbf{Une couche $5 \times 5$} :
        $5^2 = 25$ paramètres.
    \end{itemize}
    Deux couches $3 \times 3$ sont donc plus économiques ($18 < 25$) tout en couvrant
    le même champ réceptif. C'est l'une des raisons pour lesquelles les architectures
    modernes (VGG, ResNet) privilégient des petits noyaux empilés plutôt que de grands
    noyaux uniques.
  }
\end{enumerate}
}
